\section{Flagship case study: Brualdi--Goldwasser}
\label{sec:bg}

We use an in-progress formalization of the Brualdi--Goldwasser extremal-tree
conjecture as the paper's research-scale demonstration. This section states, as
plainly as we can, what has and has not been achieved. The headline is not a
completeness claim---\textbf{the conjecture is not proved}---but the fact that the
boundary between proved and open is, in this formalization, a machine-checked
object: the reader does not have to trust our prose about which lemmas are done.

\subsection{The conjecture and why it is hard}

\todo{State the Brualdi--Goldwasser conjecture precisely, with the original
citation~\cite{brualdigoldwasser}: the tree functional in question and the claimed
extremal tree. Keep this to the exact published statement---do not paraphrase into
something stronger.} At the level needed here: the conjecture asserts that a
specific graph functional, defined on all trees of a given order, is maximized by a
particular family of trees. The difficulty is that it is an \emph{extremal}
statement over an unbounded, combinatorially explosive family; direct verification
is available only for small orders (computationally, up to a bound of a few dozen
vertices), and an unconditional proof requires a structural argument that reduces
the arbitrary tree to a canonical extremal form.

Internally, the formalization works with an exact-rational tree functional
(written $\mathrm{Aobj}$) obtained by realizing a tree into a cavity/matching model.
An important and openly documented caveat lives here: several of the analytic
results below are proved about this \emph{model}, and the identification of the
model with the original combinatorial object it is meant to represent is a separate
step that is not yet formalized. We flag each such result explicitly.

\subsection{The reduction and its two open layers}

The formalization's spine is a single capstone theorem,
\code{R3Cert.Step3.conjecture1\_of\_layers}, which is \emph{conditional}: it takes
two hypotheses and derives the full conjecture from them. The theorem is
sorry-free; what it establishes is a reduction, not the conjecture. The two
hypotheses are the open layers:
\begin{description}[leftmargin=1.2em,style=nextline]
  \item[\normalfont$H_{\mathrm{norm}}$ (normalization).] Every tree is dominated,
    under the functional, by a balanced, capped ``hub-state''---a canonical form.
    Formally: for every tree $t$ there is a balanced, capped hub-state $s$ with
    $\mathrm{Aobj}(t) \le \mathrm{Aobj}(\text{backbone}(s))$.
  \item[\normalfont$H_{\mathrm{dom}}$ (rate domination).] Every merge-normal
    balanced, capped state is in turn dominated by the conjectured tie value:
    $\mathrm{Aobj} \le \mathrm{tieU}$.
\end{description}
Both remain open. $H_{\mathrm{norm}}$ decomposes further into sub-goals of mixed
status: a single-hub within-one arm-balancing argument is fully formalized and
machine-checked, but the pieces that would close the full layer---pinning the
extremal arm rates to the specific $\{4,5\}$ profile, establishing the capping
invariant, the tree-to-hub-state reduction, and multi-hub extremality---are
research-open. The predicate the codebase uses to mark the end goal,
\code{conjecture1\_proved}, is defined to be \code{False} and is asserted so by a
unit test: the formalization does not, and does not claim to, prove the conjecture.

\subsection{What Telperion certifies here}

\Telperion{} is the engine behind the concrete, certificate-bearing parts of the
reduction: exact-rational positivity floors for the per-class ledger, the per-rule
reduction certificates of the internal ``R47'' argument (spider, legs, Kelmans
corners, tiebreak, de-load), and the two analytic inequalities the reduction rests
on. Where these are inequalities over rational functions or polynomials with a
tie/equality case, they are exactly the certificate shapes of \S\ref{sec:emitters}:
P\'olya-with-zeros, box-corner bilinear reductions, exact SOS reaching the interior
tie, and finite cell subdivisions glued back together. The surrounding structural
scaffolding (the tree recursion, the cavity model, the case analysis) is
hand-written \Lean; \Telperion{} supplies the parts that are certificate
computations, which is where exact arithmetic and mechanical re-checking earn their
keep.

\subsection{The proved/open ledger}

Table~\ref{tab:ledger} is the honest summary. Every ``proved'' row is sorry-free
and, as reported in \S\ref{sec:bg-evidence}, accepted by both kernels with axiom
trace exactly $[\code{propext}, \code{Quot.sound}, \code{Classical.choice}]$.

\begin{table}[h]
\centering
\small
\begin{tabular}{@{}p{0.34\textwidth}p{0.58\textwidth}@{}}
\toprule
\textbf{Proved (kernel + \nanoda, axiom-clean)} & \textbf{Meaning} \\
\midrule
\code{R3Cert.Step3.conjecture1\_of\_layers} & The conjecture \emph{follows from}
$H_{\mathrm{norm}}$ and $H_{\mathrm{dom}}$. A reduction, conditional on the two
open layers. \\
\code{R3Cert.phi\_le\_one} ($\log\Phi \le 0$) & The $\Phi \le 1$ analytic crux
---proved about the branch-cavity \emph{model} (\todo{confirm exact statement};
model-vs-object identification deferred). \\
\code{R3Cert.CappedJointConfig.} \code{gstep\_le\_one\_achievable} & The g-step /
master-inequality crux for achievable configurations. \\
\midrule
\textbf{Open (research-open)} & \\
\midrule
$H_{\mathrm{norm}}$ (normalization layer) & $\{4,5\}$ arm-rate pinning, the capping
invariant, tree-to-hub reduction, multi-hub extremality. Single-hub within-one
balancing is done; the rest is open. \\
$H_{\mathrm{dom}}$ (rate-domination layer) & Merge-normal capped states dominated
by the tie value. Open. \\
Model-to-object identification & Identifying $\mathrm{Aobj}$ / the cavity model with
the original functional. Deferred. \\
\bottomrule
\end{tabular}
\caption{The Brualdi--Goldwasser formalization, proved vs.\ open. \code{conjecture1\_proved}
is \code{False}.}
\label{tab:ledger}
\end{table}

\subsection{Machine-verification evidence}
\label{sec:bg-evidence}

The three ``proved'' anchors are the theorems guarded by the formalization's
continuous integration, and they are re-verified through the \Comparator{}
(\S\ref{sec:indep}) on every change to the formalization. For each, \code{lean4export}
exports the proof from the built project, the axiom set is checked against the clean
whitelist, and the exported term is replayed by both \Lean's default kernel and
\nanoda. In our runs all three anchors pass: both kernels accept each, and each
reports the axiom trace $[\code{propext}, \code{Quot.sound}, \code{Classical.choice}]$
---no \code{sorryAx}, no \code{native\_decide}. The conditional capstone's dependency
cone is large but tractable: its export and two-kernel replay complete in a few
minutes of continuous-integration time. Because these statements are phrased in the
formalization's own vocabulary, the \Comparator{} runs them in self-check mode
(\S\ref{sec:indep}): the axiom whitelist and the independent second kernel are real,
independent verification; independent \emph{statement} authorship is not available
for a statement about the project's own structures, and we do not claim it.

The upshot is a precise and unfakeable claim. We do not assert the
Brualdi--Goldwasser conjecture. We assert that a specific conditional reduction and
its two analytic cruxes are proved, sorry-free, and accepted by two independent
kernels with a clean axiom trace---and that everything else, including the two
layers on which the reduction depends, is open. A separate paper will present the
result as a theorem only if and when those layers are closed.
